\documentclass[accepted]{uai2026} % camera-ready

%% Choose your variant of English; be consistent
\usepackage[american]{babel}

%% Packages
\usepackage{natbib}
    \bibliographystyle{plainnat}
    \renewcommand{\bibsection}{\subsubsection*{References}}
\usepackage{mathtools}
\usepackage{amssymb}
\usepackage{booktabs}
\usepackage{multirow}
\usepackage{graphicx}
\usepackage{tikz}
\usepackage{xcolor}
\usepackage{enumitem}
\usepackage{amsthm}
\newtheorem{theorem}{Theorem}

%% Macros
\newcommand{\ie}{i.e.\@}
\newcommand{\eg}{e.g.\@}

\title{Diagnosing Conformal Prediction Failures Under Distribution Shift:\\A COVID-19 Case Study}

% Author info hidden for anonymous submission (automatically by class)
\author[1]{Chorok Lee}
\affil[1]{%
    Korea Advanced Institute of Science and Technology (KAIST)\\
    Daejeon, South Korea
}

\affil[1]{\texttt{choroklee@kaist.ac.kr}}

\begin{document}
\maketitle

\begin{abstract}
We examine the reported association between feature-importance concentration and conformal coverage degradation in a COVID-era supply-chain case study and external datasets. Rounded historical summaries yield Spearman $\rho=0.853$ for a selected 16-task subset and $\rho=0.654$ for all 17 reported tasks. A statistical audit identified test-aware preprocessing, reuse of calibration labels for early stopping, different prediction-set conventions across cohorts, and missing numerical artifacts. These findings limit the original pre-deployment and replication claims; the historical experiments have not been rerun under a corrected common protocol. We give a conditional score bound for an idealized probability mixture, whose weight is distinct from measured SHAP concentration, and withdraw unsupported class-count and ICC interpretations. The diagnostic and 40\% threshold remain exploratory hypotheses requiring prospective validation.
\end{abstract}

\section{Introduction}\label{sec:intro}

Conformal prediction provides distribution-free coverage guarantees under the assumption of exchangeability~\citep{vovk2005algorithmic}. However, real-world deployments face distribution shifts that violate this assumption. While prior work characterizes \textit{how} conformal prediction degrades under shift~\citep{tibshirani2019conformal, barber2023conformal}, a critical gap remains: \textbf{Can we identify which deployed models will fail before observing test data?}

We propose \textbf{SHAP concentration}---the fraction of total feature importance concentrated in the top feature---as a pre-deployment diagnostic for conformal prediction vulnerability. Using a supply chain benchmark with temporal splits aligned to COVID-19 (February--July 2020), we conduct an observational case study where 8 classification tasks experience common temporal periods yet exhibit coverage drops ranging from negligible ($<$1\%) to catastrophic ($>$70\%).

Our contributions:

\begin{enumerate}
    \item \textbf{Exploratory diagnostic}: The reported SHAP concentration summaries correlate with coverage-drop summaries: $\rho=0.853$ across the selected 16 tasks and $\rho=0.654$ across all 17 tasks including Stack Overflow. These correlations can be recomputed from rounded published points, but the underlying experiment ledgers are missing and historical SALT/external set conventions differ. Prospective diagnostic value requires a common protocol without test-aware preprocessing.

    \item \textbf{Conditional theory}: Theorem~\ref{thm:score_inflation} bounds an idealized probability mixture with explicit assumptions and a weight distinct from measured SHAP concentration. It does not verify those assumptions empirically or establish a guarantee for the archived experiments. Historical shift-detector comparisons are descriptive and await reproducible outputs.

    \item \textbf{Exploratory framework}: A 40\% concentration rule and protective-factor check summarize observed cases. They are retrospective hypotheses, with one external false negative and a protective rule derived from one false positive. Historical single-seed retraining evidence is suggestive (Holm-adjusted $p=0.11$), not a validated deployment recommendation.
\end{enumerate}

\section{Related Work}\label{sec:related}

\textbf{Conformal Prediction.} \citet{vovk2005algorithmic} introduced conformal prediction with exchangeability guarantees; \citet{shafer2008tutorial} provide an accessible tutorial. Recent work extends to classification~\citep{romano2020classification,cauchois2021knowing} and regression~\citep{romano2019conformalized,lei2018distribution}. \citet{angelopoulos2023gentle} provide a comprehensive survey. \citet{ding2023class} show that class-conditional coverage guarantees become harder to achieve as class count grows and propose clustered conformal prediction as a remedy; \citet{gibbs2025conditional} provide finite-sample conditional guarantees via a covariate-shift reweighting framework interpolating between marginal and conditional validity. Both works frame class count as a challenge---a confound we explicitly examine via partial correlation analysis (Section~\ref{sec:baselines}).

\textbf{Conformal Prediction Under Shift.} \citet{tibshirani2019conformal} study covariate shift with known propensity scores. \citet{podkopaev2021distribution} develop distribution-free uncertainty quantification for classification under label shift. \citet{barber2023conformal} provide a comprehensive treatment beyond exchangeability, bounding coverage loss by the total variation distance between test and calibration score distributions. \citet{kasa2023empirically} empirically characterize how CP degrades across many vision architectures and shift types, showing substantial degradation is common---our work addresses the complementary question of which tabular models will fail \textit{before} test data is observed. We complement these theoretical characterizations with an empirical diagnostic that is associated with this distributional divergence---and hence with which tasks fail---before deployment.

\textbf{Adaptive Methods.} \citet{gibbs2021adaptive,gibbs2024conformal} develop Adaptive Conformal Inference (ACI) for non-stationary settings, with extensions by \citet{zaffran2022adaptive} for time series, \citet{feldman2023achieving} for online risk control, and \citet{bhatnagar2023improved} for strongly adaptive online learning. \citet{angelopoulos2024conformal} generalize conformal prediction to broader risk measures. \citet{kasa2025adapting} propose entropy-based methods (ECP/EACP) that adapt prediction sets using unlabeled test data at deployment time---a related but distinct goal from ours, which requires no test observations: SHAP concentration is a structural pre-deployment diagnostic computed solely on validation data. Our experiments show ACI recovers marginal coverage under severe shift but at a substantial cost in prediction set informativeness.

\textbf{Shift Detection.} WILDS~\citep{koh2021wilds} and Shifts~\citep{malinin2021shifts} provide benchmarks, while \citet{gulrajani2021search} show domain generalization methods often fail to improve over ERM under realistic distribution shifts. \citet{gardner2023tableshift} benchmark tabular distribution shift across 15 tasks---the closest tabular-specific benchmark; our SALT dataset differs in providing a naturalistic COVID-19 temporal split with documented business-disruption context rather than curated benchmark partitions. \citet{garg2022leveraging} use unlabeled test data to predict accuracy degradation---a related goal but requiring test-time observations. Kernel-based tests (MMD~\citep{gretton2012kernel}) and classifier two-sample tests (C2ST~\citep{lopez2017revisiting}) detect distributional differences but do not predict the \textit{severity} of downstream failure. We show empirically that MMD and C2ST detect shift for all tasks uniformly yet cannot distinguish catastrophic from robust conformal outcomes (Section~\ref{sec:shift_detection}). Our contribution is a \textit{pre-deployment} diagnostic using feature importance structure that predicts failure severity, not merely shift presence. \citet{miller2021accuracy} show that in-distribution (ID) validation accuracy predicts out-of-distribution (OOD) accuracy for image classifiers---a correlation that does not hold for conformal coverage in our tabular setting, where ID validation coverages cluster near the 90\% target, spanning 83.6\%--99.9\% (Table~\ref{tab:main_results}), yet test-time coverage spans 12--100\%, reinforcing the need for structure-based diagnostics beyond accuracy metrics.

\textbf{Interpretability for Reliability.} Our work connects SHAP~\citep{lundberg2017unified,lundberg2020local} to model reliability assessment, extending feature attribution from post-hoc analysis to prospective failure diagnosis. Our use is prospective: concentration is computed pre-deployment on validation data, predicting which models will fail.

\section{Methodology}\label{sec:method}

\subsection{Dataset and Temporal Splits}

We use the SALT (Supply chain ALlocaTion) dataset from RelBench~\citep{fey2024relbench}, a supply chain benchmark with temporal splits. All 8 classification tasks from the rel-salt benchmark are included; no tasks were excluded.
External analyses report 9 additional datasets. We retain the historical $K\geq4$ subset (8 external datasets plus 8 SALT tasks, $n=16$) as an exploratory scope choice. APS mechanics do not rule out catastrophic coverage failure at $K=2$ or $3$, and no dated prospective exclusion plan is available. Including Stack Overflow ($K=3$) yields the all-17 sensitivity result $\rho=0.654$, $p=0.004$; this sensitivity limits conclusions from the selected subset.
\begin{itemize}
    \item \textbf{Train}: Before February 2020 (pre-COVID)
    \item \textbf{Validation}: February--July 2020 (COVID onset)
    \item \textbf{Test}: After July 2020 (COVID peak)
\end{itemize}

\noindent\textbf{Audit of the historical data protocol.} The supplied SALT code fits categorical encoders and the target vocabulary using train, validation, and test data. It also uses all validation labels for early stopping before dividing that sample into calibration and evaluation; the external script likewise reuses validation labels. Thus, although the reported SHAP summaries use validation covariates, the full pipeline is not strictly pre-deployment and does not satisfy standard held-out split-conformal assumptions. The numerical tables below are historical reported results, not a validated rerun. Corrected evaluation requires train-only preprocessing, an independently defined class vocabulary, and separate tuning/calibration/evaluation samples. Threshold evaluation uses test-drop labels and is retrospective.

\subsection{Conformal Prediction Setup}

The archived experiments use two deterministic APS-related conventions. SALT set construction includes the class that crosses the calibrated cumulative-probability threshold. The external script instead computes coverage by the event $s(x,y)\leq\hat q$, while its size calculation includes the crossing class. These define different sets; the historical pooled endpoint and external coverage/size pairs therefore require reevaluation under one explicit convention. A new reference implementation uses the exact finite-sample order statistic and inverts a fixed-tie-order inclusive APS score. It has not generated the historical results reported here.
\begin{enumerate}
    \item Train LightGBM classifier with fixed hyperparameters
    \item Calibrate conformal predictor on 50\% of validation set
    \item Evaluate coverage on held-out validation and test sets
\end{enumerate}

\subsection{Ensemble Protocol}

The historical design reports 50 seed runs (42--91) on the same data, with a fixed chronological validation split, $t$ intervals across seeds and paired Wilcoxon comparisons. These summarize algorithmic variation conditional on one dataset and do not create 50 independent realizations of a distribution shift. Original seedwise results are unavailable for recomputing the intervals/tests.

\subsection{Feature Overlap Metric}

Feature temporal stability uses Jaccard similarity:
\begin{equation}\label{eq:jaccard}
J(x_j) = \frac{|A_{\text{train}} \cap A_{\text{test}}|}{|A_{\text{train}} \cup A_{\text{test}}|}
\end{equation}
where $A_{\text{train}}$ and $A_{\text{test}}$ are the sets of unique values for feature $x_j$ in train and test data. In the pre-deployment framework (Section~\ref{sec:framework}), we substitute validation data for test data as a proxy.

\subsection{SHAP Concentration Diagnostic}\label{sec:shap_method}

We define the SHAP concentration metric as:
\begin{equation}\label{eq:concentration}
C = \frac{\phi_1}{\sum_{j=1}^{p} \phi_j}
\end{equation}
where $\phi_j$ is the mean absolute SHAP value of the $j$-th feature computed on validation data using TreeExplainer~\citep{lundberg2020local}, and $\phi_1$ is the importance of the top-ranked feature. $C$ is computed \textbf{before deployment} using only validation data.

\subsection{Why SHAP Concentration?}\label{sec:why_shap}

A natural question is whether simpler diagnostics suffice.

\textbf{Feature overlap (Jaccard) is insufficient.} Tasks relying on transaction IDs have Jaccard $\approx 0$ while tasks with entity-based features have Jaccard $> 0.5$ (Table~\ref{tab:overlap}), yet Jaccard alone cannot explain the full range of coverage drops.

\textbf{Post-hoc metrics require test data.} Prediction entropy changes and Expected Calibration Error (ECE) shifts can detect failures---but only \textit{after} observing test data (Section~\ref{sec:baselines}). Among the diagnostics we evaluate, SHAP concentration is available before deployment and shows the strongest severity association.

\textbf{Why not raw feature importance?} Raw importance values are not comparable across tasks (different scales, feature counts). Concentration normalizes to a $[0,1]$ scale, enabling cross-task comparison.

\textbf{Why top-1?} Top-1 has the strongest observed association among the reported concentration summaries. The table reports eight candidate metrics (top-1/2/3/5/10, HHI, Gini, entropy). With the displayed SALT points, its unrounded two-sided Spearman $p=0.0101755$ gives Holm-adjusted $p=0.0814$ for this eight-test family ($0.0509$ even for the previously reported five-test family). Neither clears 0.05. The pooled $n=16$ result reuses these eight SALT tasks and is not independent replication of metric selection. On the eight external rounded summaries alone, $\rho=0.833$, $p=0.0102$; this is descriptive arithmetic, pending provenance and pipeline reconciliation.

\begin{table}[t]
\centering
\caption{Per-Task Feature Importance Concentration. Top-$k$ shows cumulative importance in the top $k$ features (\%). Drop = coverage degradation (pp). High Top-1 with catastrophic drop (e.g., s-payterms: 54.2\% $\rightarrow$ 77.1~pp) versus high Top-2/3 with near-zero drop (s-office: 77.0\%/84.8\% $\rightarrow$ 0.1~pp) is an observed contrast, not evidence of a causal mechanism.}\label{tab:topk_concentration}
\small
\begin{tabular}{@{}lcccc@{}}
\toprule
Task & Top-1 & Top-2 & Top-3 & Drop \\
\midrule
s-payterms & 54.2 & 78.1 & 89.6 & 77.1 \\
s-shipcond & 50.7 & 77.5 & 92.5 & 71.6 \\
i-shippoint & 48.8 & 72.7 & 92.4 & 18.5 \\
s-group & 47.3 & 65.9 & 79.6 & 71.2 \\
s-office & 42.6 & 77.0 & 84.8 & 0.1 \\
i-incoterms & 28.9 & 42.3 & 55.1 & 11.3 \\
i-plant & 23.9 & 46.9 & 65.0 & 10.6 \\
s-incoterms & 23.7 & 47.2 & 69.1 & 8.5 \\
\bottomrule
\end{tabular}

\raggedright
\footnotesize
Tasks sorted by Top-1 concentration. Top-$k$ = $\sum_{j=1}^{k} \phi_j / \sum_{j=1}^{p} \phi_j$ (\%). Drop = mean coverage drop (50 seeds, pp). Spearman $\rho$ with drop: Top-1 = 0.833 ($p=0.010$ before multiplicity correction); Top-2 = 0.524 ($p=0.183$); Top-3 = 0.524 ($p=0.183$); Top-5 = 0.690 ($p=0.058$); Top-10 = 0.399 ($p=0.328$). Alternative concentration metrics: HHI = 0.619 ($p=0.102$); Gini = 0.500 ($p=0.207$); Entropy = 0.571 ($p=0.139$)---none significant. The contrast for sales-office motivates a mechanism hypothesis; the stability of individual features is not established by this correlation comparison.
\end{table}

\textbf{Model scope.} The reported model-sensitivity correlations differ (LightGBM $0.833$, CatBoost $0.667$, XGBoost $0.548$, RF $0.30$, MLP $0.43$). These observational comparisons do not identify architectural causes or establish that boosting necessarily reinforces single-feature dependence, bagging necessarily dilutes it, or SHAP approximation error explains the MLP result. TreeSHAP is exact for the specified tree-explanation model, but a causal mechanism would require controlled interventions and attribution-stability checks. Generalization beyond the evaluated models remains unverified.

\section{Theory: Why Concentration Relates to Failure}\label{sec:theory}

We give a conditional score bound for an idealized probability mixture. Its weight $\lambda$ is distinct from measured SHAP concentration $C$. The calculation is not a guarantee for the archived experiments and does not establish their mechanism.

\textbf{Score and set convention.} Fix a deterministic class-index tie order. Let $s(x,y^*)$ be the sum of probabilities through the true label in decreasing-probability order, and let $\mathcal B(x,y^*)$ contain the labels ranked after it. Then
\begin{equation}\label{eq:aps_equiv}
s(x,y^*)=1-\sum_{y\in\mathcal B(x,y^*)}\hat p(y\mid x).
\end{equation}
Each omitted probability is at most $\hat p(y^*\mid x)$, including ties. The score-inversion set is $\mathcal C_q(x)=\{y:s(x,y)\leq q\}$. This convention differs from the historical SALT crossing-label set; that distinction must be resolved before using this proposition to interpret its coverage.

\begin{theorem}[Conditional bound for an idealized probability mixture]\label{thm:score_inflation}
Let $K\geq2$, $\lambda\in[0,1]$, and a fixed deterministic threshold $q\in[0,1]$. Assume
\begin{enumerate}[label=\textup{(A\arabic*)},leftmargin=*,nosep]
\item $\hat p(y\mid x)=\lambda g(y\mid x_1)+(1-\lambda)h(y\mid x_{\setminus1})$, where $g,h$ are probability distributions.
\item At test time $g(y^*\mid x_1)\leq\varepsilon$ almost surely, for fixed $\varepsilon\in[0,1]$.
\item The joint distribution of $(h(\cdot\mid x_{\setminus1}),y^*)$ is the same in calibration and test populations.
\end{enumerate}
Write $\bar h=\mathbb E[h(y^*_{\rm cal}\mid x_{\setminus1}^{\rm cal})]$. Then
\begin{equation}\label{eq:score_pointwise}
s(x^{\rm test},y^*)\geq1-(K-1)\{\lambda\varepsilon+(1-\lambda)h(y^*\mid x_{\setminus1}^{\rm test})\},
\end{equation}
and consequently
\begin{equation}\label{eq:score_bound}
\mathbb E[s_{\rm test}]\geq1-(K-1)\{\lambda\varepsilon+(1-\lambda)\bar h\}.
\end{equation}
Holding $\varepsilon$ and $\bar h$ fixed, this lower bound is strictly increasing in $\lambda$ if $\varepsilon<\bar h$. This concerns the bound, not monotonicity of the actual score expectation.

For $\lambda<1$, score-inversion coverage obeys
\begin{equation}\label{eq:cov_bound}
\mathbb P(y^*\in\mathcal C_q(x^{\rm test}))\leq
\mathbb P\!\left(h(y^*_{\rm cal})\geq T(\lambda)\right),
\end{equation}
where
\[
T(\lambda)=\frac{(1-q)/(K-1)-\lambda\varepsilon}{1-\lambda}.
\]
Holding $q$, $\varepsilon$ and the distribution of $h(y^*_{\rm cal})$ fixed, this upper bound is non-increasing when $\varepsilon\leq(1-q)/(K-1)$; the threshold is strictly increasing under strict inequality. At $q=1,\varepsilon=0$ it is constant. At $\lambda=1$, if $q<1-(K-1)\varepsilon$ coverage is zero; otherwise this argument yields only the trivial upper bound one.
\end{theorem}

\begin{proof}
The omitted sum in~\eqref{eq:aps_equiv} contains at most $K-1$ terms, each at most $\hat p(y^*)$. Substituting (A1)--(A2) gives~\eqref{eq:score_pointwise}; expectation and (A3) give~\eqref{eq:score_bound}. For the score-inversion event, $s\leq q$ implies the lower bound is at most $q$, yielding~\eqref{eq:cov_bound} when $\lambda<1$. Its threshold derivative is $[(1-q)/(K-1)-\varepsilon]/(1-\lambda)^2$. The $\lambda=1$ case follows directly from the pointwise bound.
\end{proof}

The proposition does not identify $\lambda$ with the mean-absolute TreeSHAP ratio, whose usual LightGBM decomposition is in raw-margin space. Recalibrating a changed classifier also changes $q$ and may change $h$, so the fixed-quantity monotonicity does not establish monotone empirical coverage across classifiers. The historical score comparisons in Appendix~\ref{app:theory_proof} are illustrative substitutions, not verification of (A1)--(A3).

\section{Results}\label{sec:results}

\subsection{Coverage Degradation Across Tasks}

Despite experiencing common COVID-era temporal periods, the 8 supply chain tasks exhibit dramatically different coverage degradation (Table~\ref{tab:main_results}). All coverage drops are statistically significant (paired Wilcoxon $p \leq 0.005$, 95\% CIs from the $t$-distribution), yet range from 0.1\% to 77.1\%.

\begin{table*}[t]
\centering
\caption{Coverage Degradation Under COVID-19 Distribution Shift (50 Seeds).
95\% CIs from $t$-distribution. $p$-values from paired Wilcoxon signed-rank test
(val $\rightarrow$ test, seed-level). All drops significant at $p \leq 0.005$.}\label{tab:main_results}
\small
\begin{tabular}{@{}lcccccc@{}}
\toprule
Task & Cl & Val Coverage [95\% CI] & Test Coverage [95\% CI] & Drop [95\% CI] & $p$ & Cat \\
\midrule
s-shipcond & 45 & 93.5 [93.2, 93.7] & 21.8 [14.1, 29.6] & 71.6 [63.9, 79.4] & $<$0.001 & SEV \\
s-group & 459 & 83.6 [81.7, 85.5] & 12.4 [3.1, 21.7] & 71.2 [61.3, 81.0] & $<$0.001 & SEV$^*$ \\
s-payterms & 137 & 90.8 [90.7, 91.0] & 13.7 [6.0, 21.5] & 77.1 [69.4, 84.9] & $<$0.001 & SEV \\
i-plant & 35 & 92.0 [91.7, 92.3] & 81.4 [79.0, 83.8] & 10.6 [8.2, 13.0] & $<$0.001 & ROB \\
i-shippoint & 69 & 91.2 [90.9, 91.4] & 72.7 [62.3, 83.0] & 18.5 [8.2, 28.8] & 0.005 & ROB$^{*\dagger}$ \\
s-incoterms & 13 & 95.5 [95.3, 95.7] & 87.0 [84.7, 89.3] & 8.5 [6.2, 10.8] & $<$0.001 & ROB \\
i-incoterms & 13 & 95.0 [94.9, 95.2] & 83.7 [80.9, 86.6] & 11.3 [8.5, 14.1] & $<$0.001 & ROB \\
s-office & 25 & 99.9 [99.9, 100.0] & 99.9 [99.9, 99.9] & 0.1 [0.0, 0.1] & $<$0.001 & ROB \\
\bottomrule
\end{tabular}


\raggedright
\footnotesize
SEV = Severe ($>$50\% drop), ROB = Robust ($<$20\% drop); $^*$high model variance (CV $>$ 50\%), classified by median; s-group's reported 83.6\% validation coverage is below target; high class count alone does not explain failure of a marginal conformal guarantee. Calibration reuse, chronological shift and implementation require investigation; $^\dagger$meets At-risk criterion (drop $>$15~pp) under the decision framework (Table~\ref{tab:framework_validation}).
\end{table*}

\subsection{Diagnostic Analysis}

Two factors account for much of the variance:

\textbf{Factor 1: Feature temporal stability.} Tasks using transaction IDs (Jaccard $\approx 0.02$) fail catastrophically; tasks using stable entities (products, organizations; Jaccard $> 0.5$) maintain coverage (Table~\ref{tab:overlap}).

\textbf{Factor 2: Feature importance concentration.} Among tasks with universally low Jaccard, concentration determines vulnerability (Section~\ref{sec:shap_results}).

\begin{table}[t]
\centering
\caption{Feature Overlap for Primary Features.}\label{tab:overlap}
\small
\begin{tabular}{lccc}
\toprule
Task & Primary Feature & Jaccard & Drop \\
\midrule
s-shipcond & SALESDOC (ID) & 0.02 & 71.6\% \\
s-group & SALESDOC (ID) & 0.02 & 71.2\% \\
i-incoterms & PRODUCT, PARTY & 0.58 & 11.3\% \\
s-office & SALESORG & 0.61 & 0.1\% \\
\bottomrule
\end{tabular}


\raggedright
\footnotesize
Drop = mean val $-$ mean test coverage (50 seeds).
\end{table}

\subsection{SHAP Concentration Correlates with Failure}\label{sec:shap_results}

Across all 8 SALT multiclass tasks, SHAP concentration correlates with coverage degradation: Spearman $\rho = 0.833$, $p = 0.010$ (Table~\ref{tab:stratified_correlation}, Figure~\ref{fig:n16_correlation}). Bootstrap 95\% CI: $[0.30, 1.00]$ (10,000 resamples, percentile method; BCa would yield a modestly tighter lower bound at $n=8$ but does not change the qualitative conclusion). The correlation magnitude is stable under leave-one-out analysis ($\rho \in [0.75, 0.96]$), though statistical significance is not maintained for 2 of 8 jackknife samples ($p = 0.052$) due to reduced power at $n=7$. SHAP concentration values are highly stable: bootstrap resampling (1,000 resamples of 10K validation samples) yields coefficient of variation below 1\% for all tasks, with 95\% CIs within $\pm$1 percentage point.

\textbf{Mechanism.} Figure~\ref{fig:shap} reveals the distinction. Catastrophic tasks: the dominant feature experiences a concentrated importance explosion (4.5$\times$ increase). Robust tasks: importance redistributes across features with ranking reshuffle (10--20$\times$ increases distributed). These observations do not verify the probability-mixture assumptions of Theorem~\ref{thm:score_inflation}.

\textbf{False positives.} Under the mean-drop $>$15~pp criterion, sales-office is the SALT false positive; item-shippoint is at risk by mean although most seeds have small drops. The proposed protective rule was derived from sales-office alone, and inconsistent descriptions of its feature identity and overlap split require the original feature ledger. No deployment guarantee follows from this retrospective rule.

\textbf{Dependence and uncertainty.} SALT tasks share entities, periods and one benchmark. Their statistical independence is not established by between-task ICC. Seed-level intervals describe algorithmic variability conditional on these data; they do not quantify uncertainty across new datasets or future shifts. Negative validation/test agreement ICC cannot establish negative correlation or conservative paired tests (Appendix~\ref{app:icc}).

\textbf{Inferential status.} Both the within-SALT and pooled cross-domain analyses are exploratory. The pooled endpoint includes the SALT tasks used in diagnostic development. Observed-effect power calculations do not confer confirmatory status, and no dated prespecified analysis plan or raw prediction ledger is available in the released package.

\textbf{Cross-domain extension.} The pooled historical summaries yield $\rho=0.853$ and Kendall $\tau=0.667$ for the selected 16 tasks. The supplied eight external rounded points alone yield $\rho=0.833$, $p=0.0102$, while all 17 tasks yield $\rho=0.654$, $p=0.004$. These arithmetic checks do not validate the underlying missing runs, independence of task selection, or comparability of historical APS conventions. Covertype is the reported catastrophic external case and KDDCup99 a low-concentration false negative.

\begin{table}[t]
\centering
\caption{Stratified Correlation Analysis.}\label{tab:stratified_correlation}
\resizebox{\columnwidth}{!}{%
\begin{tabular}{@{}lccccc@{}}
\toprule
Group & $n$ & $\rho$ & Kendall $\tau$ & $p$($\rho$) & Boot.\ 95\% CI \\
\midrule
Multiclass (SALT)    & 8 & 0.833 & 0.714 & 0.010 & $[0.30, 1.00]$ \\
Multiclass (4 dom.) & 11 & 0.909$^\dagger$ & 0.782 & $<$0.001 & $[0.61, 1.00]$ \\
Multiclass (8 dom.) & 15 & 0.882$^\ddagger$ & 0.714 & $<$0.001 & $[0.60, 0.97]$ \\
\textbf{Multiclass (9 dom.)} & \textbf{16} & \textbf{0.853} & \textbf{0.667} & $<$\textbf{0.001} & $\mathbf{[0.50, 0.96]}$ \\
Combined (10 dom.)$^\S$ & 17 & 0.654 & --- & 0.004 & --- \\
\bottomrule
\end{tabular}}


\raggedright
\footnotesize
LOO stability (SALT): $\rho \in [0.75, 0.96]$; 6/8 jackknife significant. Multiclass (9 domains) is the primary result, excluding Stack Overflow under the historical scope choice. Under the at-risk criterion (drop $>$15~pp), threshold performance at 40\% is Precision $= 0.83$, Recall $= 0.83$ (Appendix~\ref{app:threshold}).
$^\dagger$Single-seed external values; multi-seed-consistent value at $n=11$ is $\rho=0.818$, $p=0.002$ (see Appendix~\ref{app:icc}).
$^\ddagger$$n=15$ is an intermediate analysis (7 external multiclass domains); $n=16$ adds KDDCup99 as the 8th external domain to complete the 9-domain primary endpoint. Stack~Overflow (3 classes, excluded scope subset) is excluded from all multiclass endpoints.
$^\S$Combined row includes all 17 tasks (8 SALT + 9 external, 10 domains: SALT $+$ 9 external), adding Stack~Overflow back to the $n=16$ primary endpoint. Spearman $\rho=0.654$, $p=0.004$ (from Appendix~\ref{app:icc}); Kendall $\tau$ not reported for this auxiliary endpoint.
\end{table}

\begin{figure*}[t]
\centering
\includegraphics[width=0.85\textwidth]{figures/figure3_feature_importance}
\caption{\textbf{Feature Importance Dynamics Under Shift.} (A) Catastrophic task: dominant feature explodes 4.5$\times$ while maintaining top rank (concentrated dependence). (B) Robust task: importance redistributes across features with ranking reshuffle (diversified dependence). (C--D) Importance dynamics and ranking stability. Despite $\sim$0\% Jaccard similarity, coverage drops range from negligible to catastrophic.}\label{fig:shap}
\end{figure*}

\begin{figure}[t]
\centering
\includegraphics[width=\linewidth]{figures/figure_n16_correlation}
\caption{\textbf{SHAP Concentration Correlates with Coverage Degradation (Primary Result).} Spearman $\rho=0.853$, $p<0.001$ across all 16 multiclass tasks in 9 domains (dark circles: 8 SALT supply-chain tasks; orange triangles: 8 external domains, 10-seed means). The 40\% threshold (dashed) is an exploratory historical cutoff. Within SALT alone: $\rho=0.833$, $p=0.010$ ($n=8$).}\label{fig:n16_correlation}
\end{figure}

\subsection{Diagnostic Comparison}\label{sec:baselines}

We compare SHAP concentration against two alternative pre-deployment diagnostics that require no test data (Table~\ref{tab:diagnostic_comparison}).

\textbf{Native feature importance concentration}: top-1 LightGBM gain divided by total gain (zero SHAP cost, single seed; seed=42). Achieves $\rho = 0.667$ ($p = 0.071$)---directionally correct but not significant at $\alpha = 0.05$. The weaker discrimination arises because native gain conflates feature cardinality with predictive importance: sales-office shows 66.1\% native concentration (vs.\ 42.6\% SHAP) despite near-zero coverage drop.

\textbf{Ensemble disagreement}: standard deviation of validation coverage across 50 seeds. Achieves $\rho = 0.452$ ($p = 0.260$). Disagreement captures model instability but not the feature-level mechanism driving conformal failure.

\textbf{Class cardinality as a possible confound.} The historical within-SALT partial correlations are non-significant ($0.629$, $p=0.131$ for concentration; $0.334$, $p=0.464$ for $\log K$). The reported pooled partial correlation is $0.771$ ($p=0.0008$) for concentration and $-0.010$ ($p=0.97$) for $\log K$. These are conditional associations, not evidence that concentration causally drives severity independently of model, domain, shift magnitude, or pipeline differences. Their numerical provenance and uncertainty require the original task/seed ledger.

\begin{table}[t]
\centering
\caption{Pre-Deployment Diagnostic Comparison ($n=8$). Spearman $\rho$ with 50-seed mean coverage drop. $\rho_{\text{partial}}$: partial Spearman controlling for the other variable in the SHAP--class-count pair.}\label{tab:diagnostic_comparison}
\small
\begin{tabular}{@{}lcccc@{}}
\toprule
Diagnostic & $\rho$ & $p$ & $\rho_{\text{partial}}$ & $p_{\text{partial}}$ \\
\midrule
SHAP concentration & \textbf{0.833} & \textbf{0.010} & 0.629 & 0.131 \\
$\log(\text{num\_classes})$ & 0.743 & 0.035 & 0.334 & 0.464 \\
Native FI concentration & 0.667 & 0.071 & --- & --- \\
Ensemble disagreement & 0.452 & 0.260 & --- & --- \\
\bottomrule
\end{tabular}
\end{table}

We additionally compare against two \textit{post-hoc} baselines computed from 10-seed ACI experiments:

\textbf{Prediction entropy change}: $\Delta H = H_{\text{test}} - H_{\text{val}}$, where $H = -\sum_c p_c \log p_c$ averaged over instances.

\textbf{Expected Calibration Error change}: \mbox{$\Delta \text{ECE} = \text{ECE}_{\text{test}} - \text{ECE}_{\text{val}}$}.

Both metrics can detect shift \textit{after} test data is observed, but neither is available at deployment time. SHAP concentration is computed solely on validation data, making it actionable for pre-deployment risk assessment. Comparison across 7 tasks in Appendix~\ref{app:baselines}.

\section{Extended Experiments}\label{sec:extended}

\subsection{Adaptive Conformal Inference Across All Tasks}

The historical clipped ACI variant reports coverage improvements on severe tasks with larger sets. It clips $\alpha_t$ to $[0.001,0.999]$ and quantile levels to $[0,1]$, so the original arbitrary-shift ACI guarantee does not automatically transfer. It assumes immediate labels in the processed test order. Missing prediction ledgers prevent rechecking the reported tradeoff, and a deployability judgment requires task-specific utility and feedback timing.

\subsection{Alternative Conformal Scoring}\label{sec:raps}

Regularized Adaptive Prediction Sets (RAPS)~\citep{angelopoulos2021uncertainty} shows mechanism-dependent benefit: substantially reduces drops for high-cardinality tasks (s-group: $73.5\%\rightarrow10.4\%$; s-payterms: $79.9\%\rightarrow35.1\%$) but not the most concentrated task (s-shipcond: $60.4\%\rightarrow67.8\%$). This matches Section~\ref{sec:theory}: regularization controls class-accumulation effects, not concentrated single-feature dependence (Appendix~\ref{app:raps}; RAPS used random 50/50 calibration splits vs.\ deterministic main-experiment splits---within-experiment comparisons are unaffected).

\subsection{Comparison with Standard Shift Detection}\label{sec:shift_detection}

The historical summary reports MMD $p<0.002$, C2ST accuracy $>99.9\%$, PSI $>0.002$, and weak severity correlations on eight SALT tasks. Raw taskwise detector outputs are unavailable and the supplied detector script does not parse. A nonzero PSI alone is not a significance test. These observations, if reproduced, would describe this suite; they do not establish that generic shift detectors cannot predict severity in other settings.

\subsection{Retraining Analysis}

Quarterly retraining improves s-shipcond by +18.9~pp ($p=0.036$, unadjusted, single-seed; Holm-corrected $p=0.11$ over 3 severe tasks), partially improves s-payterms (+15.5~pp, non-significant), and is ineffective for s-group ($0.0\%\rightarrow2.2\%$). Improvement concentrates in the final two months of the window; robust tasks remain near-ceiling without intervention.

\subsection{Cross-Domain and External Validation}\label{sec:cross_domain}

Historical external reports cover nine datasets, with the selected eight multiclass datasets contributing to the pooled association. Covertype and Gas Sensor have supplied geographic/temporal loader definitions. The supplied KDD loader uses a positional split of one file rather than the claimed official train/test separation; PAMAP2 and the four reported null-shift controls lack generating code/results in this package. Whether those controls experience no relevant shift cannot be established from split names alone. These gaps prevent interpreting the pooled result as a validated cross-domain replication.

\subsection{Placebo Test}

The historical pre-COVID placebo reports lower degradation for seven tasks, with ratios about 6--143 where the rounded denominator is positive; sales-office has a reported ratio of one. These summaries and asymmetric training windows do not isolate a causal COVID effect, and the generating code/results are unavailable.

\section{Decision Framework}\label{sec:framework}

The following rule is a retrospective hypothesis requiring validation before operational use. Its historical evaluation depends on the unresolved implementation/provenance issues above:

\begin{enumerate}
    \item \textbf{Compute} $C = \phi_1 / \sum_j \phi_j$ (Eq.~\ref{eq:concentration}). If $C \leq 40\%$: lower observed risk. If $C > 40\%$: check protective factors---secondary features with train-validation Jaccard $> 0.5$ and importance $> 15\%$ (thresholds empirically derived from a single observed case; treat as provisional); absent these, flag as \textbf{vulnerable}. If $C \in [35\%, 45\%]$: treat as \textbf{uncertain risk} with monitoring regardless of other factors.
    \item \textbf{If vulnerable}: consider quarterly retraining (suggestive effect; Holm-corrected $p=0.11$; single-seed; Appendix~\ref{app:framework}).
    \item \textbf{Monitor} empirical coverage continuously; defer retraining for lower-risk tasks.
\end{enumerate}

\section{Conclusion}\label{sec:conclusion}

The archived results suggest an association between concentrated feature importance and coverage degradation in the evaluated settings. The selected 16-task association is $\rho=0.853$, compared with $0.654$ when the excluded three-class task is retained. These are reproducible summary arithmetic, not an independently reproduced common-protocol experiment. Test-aware preprocessing, calibration reuse, differing prediction-set conventions and missing numerical artifacts remain material limitations. The idealized theorem concerns a separate mixture weight under explicit assumptions and does not establish the proposed empirical mechanism.
Controlled interventions, a leakage-free common pipeline, and independent prospective datasets are needed to assess mechanism, model scope and operational thresholds. The current revision corrects statistical interpretation while preserving historical results as unverified reports.

\begin{contributions}
    C.L. conceived the research question, designed and executed all experiments, performed the statistical analysis, and wrote the paper.
\end{contributions}

\begin{acknowledgements}
    We thank the RelBench team for providing the SALT dataset and the anonymous reviewers for their constructive feedback. The released supplementary package is partial; missing artifacts and detected implementation issues are documented in Appendix~\ref{app:reproducibility}.
\end{acknowledgements}

% References
\bibliography{references}

\newpage
\onecolumn

\title{Diagnosing Conformal Prediction Failures Under Distribution Shift:\\A COVID-19 Case Study (Supplementary Material)}
\maketitle

\appendix

\section{Reproducibility Details}\label{app:reproducibility}

The released supplement is incomplete. It contains six Python scripts but no original prediction ledger, checkpoints, split manifests or numerical result files. The packaged shift-detector script has a syntax error, the SHAP script depends on an absent generator, and only four of the nine external loaders are supplied. The model-sensitivity, WILDS, retraining, placebo and several ablation analyses cannot currently be reproduced from this package. The audit and synthetic regression checks are included separately; they do not reproduce the historical experiments.

\subsection{LightGBM Hyperparameters}

All models trained with fixed LightGBM~\citep{ke2017lightgbm} settings: objective=multiclass, boosting=gbdt, num\_leaves=31, learning\_rate=0.05, feature\_fraction=0.8, bagging\_fraction=0.8, bagging\_freq=5, num\_boost\_round=500, early\_stopping\_rounds=50, seeds 42--91 (50 seeds). No task-specific tuning was performed.

\subsection{Conformal Prediction}

Historical calibration uses the first half of validation records and evaluation the second half, after the same validation labels have already been used for early stopping. This chronological split need not be exchangeable even within validation. Historical code uses NumPy interpolated quantiles at $\min(\lceil(n+1)(1-\alpha)\rceil/n,1)$. The corrected reference instead selects the exact $\lceil(n+1)(1-\alpha)\rceil$th order statistic, using $+\infty$ when the required rank exceeds $n$; it also applies one score-inversion convention to coverage and set size. No historical results were replaced by this reference.

\noindent\textbf{Calibration and test set sizes.} Sales-level tasks (s-shipcond, s-payterms, s-group, s-incoterms, s-office): $n_{\text{val}} = 71{,}474$, $n_{\text{cal}} = 35{,}737$, $n_{\text{test}} = 88{,}942$. Item-level tasks (i-plant, i-shippoint, i-incoterms): $n_{\text{val}} = 293{,}896$, $n_{\text{cal}} = 146{,}948$, $n_{\text{test}} = 402{,}855$. At these calibration set sizes, the finite-sample correction is negligible: $q - (1-\alpha) < 3 \times 10^{-5}$ for all tasks, so finite-sample effects do not contribute to the observed coverage gaps.

\subsection{SHAP Computation}

Method: TreeExplainer~\citep{lundberg2020local}. Subsample: 10,000 per split. Aggregation: mean absolute SHAP across classes. Concentration metric (Eq.~\ref{eq:concentration}): computed on validation data only.

\subsection{Baseline Diagnostics}

Native feature importance concentration uses a single model (seed=42) with LightGBM's gain-based importance; concentration is top-1 gain divided by total gain. Ensemble disagreement is the standard deviation of validation coverage across 50 seeds. Both are computed without SHAP.

\subsection{Computational Resources}

Standard CPU (8 cores, 8GB RAM). Wall-clock: $\sim$3--4 hours for full 8-task, 50-seed suite. Software: Python 3.9, LightGBM 3.3, SHAP 0.41.

\subsection{External Dataset Protocols}\label{app:external_data}

Table~\ref{tab:external_datasets} retains the reported external dataset protocols. The supplied code does not establish a common pipeline: SALT and external coverage use different set conventions, validation labels are reused for tuning, and most external loaders/results are absent. In particular, the supplied KDD loader contradicts the previously stated official-test protocol, as noted below. Historical results need provenance reconciliation and a common-protocol rerun.

\begin{table}[htbp]
\centering
\caption{External Dataset Protocols. Historical reported protocols; supplied external code halves validation for calibration/evaluation; $n_\text{tr}$, $n_\text{te}$ = approximate train/test sizes used; Seeds = 10 for all external datasets. Type: \textbf{DS} = documented real-world shift; \textbf{NC} = null-shift control (random/pre-defined split, no genuine shift; low concentration and robust coverage are the expected and observed outcome).}\label{tab:external_datasets}
\small
\begin{tabular}{@{}lp{4.8cm}clllc@{}}
\toprule
Dataset & Shift mechanism & Type & $K$ & $n_\text{tr}$ & $n_\text{te}$ & Source \\
\midrule
Covertype       & Geographic: areas 1--2 (train), area 3 (val/cal), area 4 (test) & DS & 7 & 290,680 & 36,968 & \citep{dua2017uci} \\
Gas Sensor      & Temporal batch drift: batches 1--6 (train), 7--8 (cal), 9--10 (test) & DS & 6 & 9,840 & 4,070 & \citep{dua2017uci} \\
KDDCup99        & Supplied code: positional 60/20/20 of one 10\% file; shift unverified & --- & 11 & 395,073 & 98,769 & \citep{dua2017uci} \\
PAMAP2          & Population shift: subject-based split (unseen subjects at test) & DS & 12 & 1,674,379 & 268,493 & \citep{dua2017uci} \\
Shuttle         & Null-shift control: UCI Statlog; 60/20/20 random split & NC & 7 & 46,400 & 11,600 & \citep{dua2017uci} \\
Avila Bible     & Null-shift control: copyist identification; 50/50 random split & NC & 12 & 10,430 & 10,437 & \citep{dua2017uci} \\
Pendigits       & Null-shift control: UCI pre-defined split (digit recognition) & NC & 10 & 7,494 & 3,498 & \citep{dua2017uci} \\
Satimage        & Null-shift control: satellite image; 60/20/20 random split & NC & 6 & 5,144 & 1,286 & \citep{dua2017uci} \\
Stack Overflow  & Category shift (3-class subset); excluded from $n=16$ historical selected endpoint & --- & 3 & 65,105 & 127,370 & \citep{dua2017uci} \\
\bottomrule
\end{tabular}


\raggedright
\footnotesize
For Covertype, the supplied loader uses wilderness areas 1--2 for training, 3 for validation and 4 for testing. These geographic indicators are not the cover-type labels (e.g., Spruce-Fir or Lodgepole Pine). For KDDCup99, the supplied loader fetches one 10\% dataset, filters classes with at least 100 observations over the whole file, encodes object columns before splitting, and uses positional 60/20/20 partitions; no separate official test file is fetched. Historical table sizes and results need reconciliation with run provenance. The remaining reported splits are retained as historical descriptions where loaders are absent; SHAP uses the validation pool, which the external code also uses for early stopping before calibration/evaluation.
\end{table}

\section{Placebo Test Details}\label{app:placebo}

\begin{table}[htbp]
\centering
\caption{Placebo Test: Pre-COVID vs COVID Degradation (50-Seed Mean Drops).}\label{tab:placebo}
\begin{tabular}{lccc}
\toprule
Task & Placebo & COVID & Ratio \\
\midrule
s-shipcond & 0.5\% & 71.6\% & 143$\times$ \\
s-group & 2.0\% & 71.2\% & 36$\times$ \\
s-payterms & 0.0\% & 77.1\% & $>$100$\times$ \\
i-plant & 1.8\% & 10.6\% & 6$\times$ \\
i-shippoint & 1.5\% & 18.5\% & 12$\times$ \\
s-incoterms & 1.2\% & 8.5\% & 7$\times$ \\
i-incoterms & 0.6\% & 11.3\% & 19$\times$ \\
s-office & 0.1\% & 0.1\% & 1$\times$ \\
\bottomrule
\end{tabular}
\end{table}

\section{High Model Variance Analysis}\label{app:variance}

Tasks marked with $^*$ in Table~\ref{tab:main_results} exhibit extreme model variance (CV $>$ 50\%):
\begin{itemize}
    \item \textbf{s-group}: Mean 12.4\%, median \textbf{0.5\%}. Most models fail; outliers inflate mean.
    \item \textbf{i-shippoint}: Mean 72.7\%, median \textbf{89.9\%}. Most models succeed; failures deflate mean.
\end{itemize}

\section{Threshold Sensitivity Analysis}\label{app:threshold}

Table~\ref{tab:threshold_sensitivity} shows threshold performance across $n=16$ multiclass tasks. Note that the 40\% threshold is derived from the SALT concentration gap (in-sample); the 8 external datasets are fully out-of-sample. On the external subset alone at 40\%: Precision~$=$~1.00, Recall~$=$~0.50 (1 TP: Covertype; 1 FN: KDDCup99; 6 TN). Low external recall reflects sparse catastrophic cases; the primary claim rests on the $n=16$ correlation ($\rho=0.853$) rather than threshold classification. Ground-truth labels use the at-risk rule (mean coverage drop $>$15~pp, consistent with 50-seed means in Table~\ref{tab:main_results}). For high-variance tasks, the mean and median labels can differ: i-shippoint has mean drop 18.5~pp (At-risk) but median drop 1.2~pp (robust); Table~\ref{tab:main_results} labels it ROB by median for descriptive purposes, while the framework uses the mean-based threshold for threshold evaluation.

\begin{table}[htbp]
\centering
\caption{Threshold Sensitivity for SHAP Concentration Across 16 Multiclass Tasks (8 SALT + 8 external; Step 2 only, no protective-factor check). At-risk defined as coverage drop $>$15\%.}\label{tab:threshold_sensitivity}
\small
\begin{tabular}{lcccccc}
\toprule
Threshold & Prec & Rec & F1 & TP & FP & FN \\
\midrule
25\% & 0.63 & 0.83 & 0.71 & 5 & 3 & 1 \\
30\% & 0.71 & 0.83 & 0.77 & 5 & 2 & 1 \\
35\% & 0.83 & 0.83 & 0.83 & 5 & 1 & 1 \\
\textbf{40\%} & \textbf{0.83} & \textbf{0.83} & \textbf{0.83} & \textbf{5} & \textbf{1} & \textbf{1} \\
45\% & 1.00 & 0.83 & 0.91 & 5 & 0 & 1 \\
50\% & 1.00 & 0.33 & 0.50 & 2 & 0 & 4 \\
\bottomrule
\end{tabular}


\raggedright
\footnotesize
At 40\%, the single FP is s-office (concentration 42.6\%) which has a protective secondary feature (SALESORG Jaccard $= 0.61$), and the single FN is KDDCup99 (intermediate regime; $C=21.1\%$, drop $=15.9$~pp). At 30\%, the two FPs are s-office and Shuttle ($C=30.7\%$); at 25\%, i-incoterms ($C=28.9\%$) is additionally flagged. Applying the protective-factor check at 40\% removes s-office (precision $= 1.00$) but leaves the KDDCup99 FN (recall $= 0.83$).
\end{table}

\subsection{Leave-One-Out Cross-Validation}\label{app:loocv}

To assess whether the 40\% threshold generalizes beyond post-hoc selection, we perform leave-one-out cross-validation on the SALT dataset ($n=8$). For each fold, we exclude one task, compute the optimal threshold from the remaining 7 tasks, and predict the held-out task's vulnerability status.

\textbf{Summary.} LOO-CV achieves 87.5\% accuracy (7/8 correct predictions). The threshold remains stable across folds: $43.1\% \pm 5.0\%$ (mean $\pm$ std). The effect size separating failed from robust tasks is Cohen's $d = 3.08$ (large), with bootstrap 95\% CI $[0.31, 1.00]$ for the underlying correlation.

\begin{table}[htbp]
\centering
\caption{LOO-CV Per-Fold Predictions (SALT, $n=8$). Threshold = optimal threshold from 7 training tasks; Pred = predicted status at that threshold; Actual = observed outcome (Failed if drop $>$15~pp).}\label{tab:loocv}
\small
\begin{tabular}{@{}lcccccc@{}}
\toprule
Task & Conc. (\%) & Drop (pp) & Thresh (\%) & Pred & Actual & Correct \\
\midrule
s-shipcond & 50.7 & 71.6 & 45 & Vuln & Failed & \checkmark \\
s-group & 47.3 & 71.2 & 45 & Vuln & Failed & \checkmark \\
s-payterms & 54.2 & 77.1 & 45 & Vuln & Failed & \checkmark \\
i-plant & 23.9 & 10.6 & 45 & Robust & Robust & \checkmark \\
i-shippoint & 48.8 & 18.5 & 45 & Vuln & Failed & \checkmark \\
s-incoterms & 23.7 & 8.5 & 45 & Robust & Robust & \checkmark \\
i-incoterms & 28.9 & 11.3 & 45 & Robust & Robust & \checkmark \\
\textbf{s-office} & \textbf{42.6} & \textbf{0.1} & \textbf{30} & \textbf{Vuln} & \textbf{Robust} & $\times$ \\
\bottomrule
\end{tabular}
\end{table}

\textbf{Effect size analysis.} Failed tasks ($n=4$) have mean concentration $50.2\% \pm 2.6\%$, while robust tasks ($n=4$) have mean concentration $29.8\% \pm 7.7\%$. The separation of 20.4~pp yields Cohen's $d = 3.08$ (large effect), describing separation in these retrospectively labeled eight tasks; it does not establish prospective reliability.

\textbf{Misclassification analysis.} Sales-office is a false positive under the concentration-only rule. Its stable-feature explanation is a hypothesis derived from this single case. The historical text inconsistently describes SALESORGANIZATION as dominant or secondary and its Jaccard overlap as train-test or train-validation; the original feature ledger is required to resolve this. A deployable protective-factor rule must use information available before test outcomes. High concentration is not necessary for failure, as the reported low-concentration KDD false negative demonstrates.

\section{Decision Framework Validation}\label{app:framework}

\begin{table}[htbp]
\centering
\caption{Concentration Check Applied to All Multiclass Tasks (Actual = At-risk if drop $>$15~pp).}\label{tab:framework_validation}
\small
\begin{tabular}{@{}lcccc@{}}
\toprule
Task & Conc. (\%) & Step 2 & Actual ($>$15~pp) & Note \\
\midrule
\multicolumn{5}{l}{\textit{SALT (in-sample)}} \\
s-payterms & 54.2 & VULN & At-risk & --- \\
s-shipcond & 50.7 & VULN & At-risk & --- \\
i-shippoint & 48.8 & VULN & At-risk$^*$ & High variance$^\dagger$ \\
s-group & 47.3 & VULN & At-risk$^*$ & --- \\
s-office & 42.6 & VULN & ROB & Protective$^\ddagger$ \\
i-incoterms & 28.9 & ROB & ROB & --- \\
i-plant & 23.9 & ROB & ROB & --- \\
s-incoterms & 23.7 & ROB & ROB & --- \\
\midrule
\multicolumn{5}{l}{\textit{External (held-out)}} \\
Covertype & 49.8 & VULN & At-risk & 10/10 seeds \\
Shuttle & 30.7 & ROB & ROB & 9/10 seeds \\
Avila Bible & 20.5 & ROB & ROB & 10/10 seeds \\
PAMAP2 & 19.2 & ROB & ROB & 10/10 seeds \\
KDDCup99 & 21.1 & ROB & At-risk$^\S$ & Seed-dependent (FN) \\
Pendigits & 14.5 & ROB & ROB & 10/10 seeds \\
Satimage & 9.0 & ROB & ROB & 10/10 seeds \\
Gas Sensor & 7.3 & ROB & ROB & 10/10 seeds \\
Stack Overflow & 48.9 & VULN & ROB & Excluded-scope false positive$^\P$ \\
\bottomrule
\end{tabular}


\raggedright
\footnotesize
$^\dagger$Median drop 1.2\%; most seeds maintain coverage.
$^\ddagger$SALESORGANIZATION: train-validation Jaccard$=$0.61, importance$>$15\%.
$^\S$Concentration below threshold ($C=21.1\pm7.5\%$), but mean drop exceeds the at-risk cutoff (15.9~pp); outcome seed-dependent (range $-$0.8 to 73.5~pp).
In-sample F1 $= 0.89$ (1 FP: s-office). External: Covertype deterministic catastrophic (10/10); five low-concentration datasets deterministic robust (10/10); Shuttle near-deterministic (9/10); KDDCup99 is the intermediate-regime false negative.
$^\P$Stack Overflow ($K=3$, $C=48.9\%$) is a reported false positive with drop $-7.0$~pp. Its exclusion from the historical $n=16$ subset is not justified by a general APS ceiling: binary and three-class predictors can also lose all coverage under shift. The all-17 association is $\rho=0.654$, $p=0.004$.
\end{table}

\noindent\textbf{Seed stability protocol.} KDDCup99's multi-seed analysis ($C=21.1\pm7.5\%$) reveals that single-seed SHAP concentration can be unstable. We recommend computing $C$ over $k \geq 5$ seeds and using the mean. Tasks in the ambiguous zone (35--45\%) should receive additional scrutiny: if the upper 95\% CI exceeds 40\%, treat as potentially vulnerable. This safety override is designed to reduce KDDCup99-like false negatives in intermediate regimes. For SALT tasks, concentration is stable across 50 seeds (within-task CV $<$ 5\%); instability arises primarily when the top-feature identity varies across seeds.

\noindent\textbf{Threshold selection.} The 40\% threshold is proposed based on the empirical distribution of SHAP concentration values on validation data: a natural gap separates low-concentration tasks (24--29\%) from high-concentration tasks (43--54\%), visible in Table~\ref{tab:framework_validation}. We emphasize that this threshold is \textit{exploratory}, derived from $n=8$ multiclass tasks, and should be validated before deployment in new domains.

\noindent\textbf{Cross-domain transfer test.} The 40\% threshold applied without tuning to 9 external domains yields 7/9 deterministic or near-deterministic outcomes (6 at 10/10, Shuttle at 9/10), with Covertype correctly flagged as catastrophic and Gas Sensor correctly unflagged as robust. KDDCup99 is the principal intermediate-regime false negative, while Stack Overflow (3 classes) is a false positive excluded from the selected endpoint. Across the $n=16$ multiclass set: $\rho = 0.853$, $p < 0.001$ (Section~\ref{sec:cross_domain}). The sensitivity analysis in Appendix~\ref{app:threshold} reports threshold performance across 30--50\%; given sparse catastrophic external cases, this threshold should be treated as provisional.

\section{Pre-Deployment vs Post-Hoc Baselines}\label{app:baselines}

SHAP concentration is computed on validation data \textit{before} deployment. Prediction entropy and ECE require observing test data. This section reports entropy and ECE changes for completeness, noting that they are not available as pre-deployment diagnostics.

\begin{table}[htbp]
\centering
\caption{Pre-Deployment (SHAP) vs Post-Hoc (Entropy, ECE) Shift Diagnostics (10-Seed ACI Experiments). SHAP concentration is computed on validation data only. $\Delta$Entropy and $\Delta$ECE require test-time observations.}\label{tab:baselines}
\small
\begin{tabular}{@{}lcccc@{}}
\toprule
Task & Cat & SHAP Conc.$^\dagger$ & $\Delta$Entropy & $\Delta$ECE \\
\midrule
\multicolumn{5}{l}{\textit{Severe tasks}} \\
s-shipcond & SEV & 50.7\% & $-$0.746 & +0.408 \\
s-group & SEV & 47.3\% & $-$1.274 & +0.176 \\
s-payterms & SEV & 54.2\% & $-$1.979 & +0.629 \\
\midrule
\multicolumn{5}{l}{\textit{Robust tasks}} \\
i-plant & ROB & 23.9\% & +0.181 & +0.131 \\
i-shippoint & ROB & 48.8\% & +0.469 & +0.073 \\
s-incoterms & ROB & 23.7\% & +0.364 & +0.350 \\
s-office & ROB & 42.6\% & +0.032 & +0.104 \\
\bottomrule
\end{tabular}


\raggedright
\footnotesize
$^\dagger$Pre-deployment diagnostic (computed on validation data only). item-incoterms omitted (3-seed pilot only).
\end{table}

Counter-intuitively, all three catastrophic tasks (coverage drop $>$ 70\%) show \textit{decreasing} prediction entropy ($\Delta H < 0$): the model becomes \textit{more} confident at test time, not less. This occurs because concentrated dependence on an OOD feature produces arbitrary but confident predictions. In contrast, the two robust tasks with non-trivial drops (i-plant, i-shippoint; coverage drops of 10--19\%) show \textit{increasing} entropy ($\Delta H > 0$), consistent with standard distributional broadening. A practitioner monitoring only entropy would correctly detect shift in moderate cases but would be \textit{misled} in the most damaging catastrophic cases---observing decreasing uncertainty while coverage collapses. ECE increases across all tasks but does not discriminate between severity levels (both catastrophic and robust tasks show increases of +0.07 to +0.63). In this SALT comparison, SHAP concentration is the only evaluated pre-deployment metric that flags all three catastrophic tasks.

\section{Task Independence Analysis}\label{app:icc}

We analyze pseudo-replication in the 50-seed design using intraclass correlation (ICC).

\textbf{Between-task ICC.} The archived summary reports ICC $=0.675$ for coverage drops and $0.720$ for test coverage. These describe a fitted variance decomposition across task/seed levels; they do not prove that tasks sharing one benchmark are independent. A design-effect conversion of 400 seed/task measurements cannot increase the number of independently observed tasks beyond the eight available tasks.

\textbf{Per-task ICC(val,test).} The previous analysis interpreted negative one-way agreement ICC as negative seedwise correlation and converted 50 pairs into effective sample sizes of 55--250. That inference is invalid. For example, when test coverage equals validation coverage minus a constant, Pearson correlation is one while the supplied agreement ICC can be strongly negative. We therefore withdraw the effective-sample-size table and the assertion that paired Wilcoxon tests must be conservative. Appropriate paired inference uses the distribution of seedwise differences and remains conditional on the common data.



\textbf{Multiple comparisons.} The reported concentration search includes eight metrics. Using the unrounded SALT $p$-value, the smallest Holm-adjusted value is $8\times0.0101755=0.0814$. The original five-metric family would give $0.0509$, not a value below 0.05. Intermediate and pooled subset analyses are retained as exploratory; reusing discovery tasks and selecting endpoints cannot be repaired solely by correcting three subset $p$-values.

\textbf{Multi-seed aggregation.} The primary $n=16$ correlation ($\rho = 0.853$) uses 10-seed means for all external datasets and 50-seed means for SALT tasks. The early $n=11$ result ($\rho = 0.909$) used single-seed external values; switching to multi-seed means at $n=11$ yields $\rho = 0.818$, $p = 0.002$, with the reduction driven by KDDCup99 whose multi-seed mean drop (15.9~pp) exceeds its single-seed value ($-$0.8~pp). This intermediate-concentration regime ($C=21.1\pm7.5\%$) exhibits high outcome variance characteristic of the boundary between robust and vulnerable. Extending to $n=16$ with consistent multi-seed means yields $\rho = 0.853$ ($p < 0.001$). Including Stack Overflow (3 classes, excluded scope subset) to make $n=17$ weakens correlation to $\rho = 0.654$.

\section{Theory Details and Score CDFs}\label{app:theory_proof}

\textbf{Scope of the proof.} The proof in Section~\ref{sec:theory} uses a fixed deterministic rank order and a probability-mixture weight $\lambda$. It does not assume that classes tied with the true label are strictly lower in probability, and it handles $\lambda=1$ separately. Its coverage statement is for the explicit score-inversion set with fixed $q$, not a claim about recalibrated models or the archived crossing-label implementation.

\textbf{Illustrative bound calculation, not empirical verification.} The historical comparison substituted mixture weight $\lambda=C$, $\varepsilon=0$ and $\bar h=1/K$, obtaining numbers below five reported mean scores. None of these substitutions verifies the probability-mixture assumptions. In particular the lower bound decreases with $\bar h$, so $1/K$ is not conservative for every residual model whose true-class probability is at least uniform. These comparisons do not validate a theorem for measured TreeSHAP concentration; empirical verification requires an identified residual decomposition and a consistent score/set implementation.

\begin{figure}[htbp]
\centering
\includegraphics[width=\linewidth]{figures/conformity_score_cdfs}
\caption{\textbf{Conformity score CDFs (seed=42).} Catastrophic tasks show rightward test-score shifts (stochastic dominance), while robust tasks can show near-degenerate or left-shift behavior.}\label{fig:cdfs}
\end{figure}

\section{Conformity Score Stochastic Dominance}\label{app:ks}

Table~\ref{tab:ks} reports two-sample KS statistics comparing calibration and test APS conformity scores (seed=42, non-randomized APS). Higher conformity scores indicate the model requires more cumulative probability mass to reach the true label---i.e., the model is less confident about the correct class.

\begin{table}[htbp]
\centering
\caption{KS Tests: Calibration vs.\ Test Conformity Score Distributions (Seed=42).}\label{tab:ks}
\small
\begin{tabular}{@{}lcccc@{}}
\toprule
Task & KS & $p$ & Mean$_{\text{cal}}$ & Mean$_{\text{test}}$ \\
\midrule
\multicolumn{5}{l}{\textit{Catastrophic tasks (test scores $>$ calibration)}} \\
s-shipcond & 0.956 & $<10^{-10}$ & 0.54 & 0.98 \\
s-payterms & 0.748 & $<10^{-10}$ & 0.65 & 0.92 \\
s-group & 0.676 & $<10^{-10}$ & 0.98 & 1.00 \\
\midrule
\multicolumn{5}{l}{\textit{Robust tasks}} \\
i-plant & 0.741 & $<10^{-10}$ & 0.73 & 0.86 \\
s-incoterms & 0.633 & $<10^{-10}$ & 0.75 & 0.87 \\
s-office & 0.994 & $<10^{-10}$ & 1.00 & 1.00 \\
i-incoterms & 0.559 & $<10^{-10}$ & 0.78 & 0.66 \\
i-shippoint & 0.505 & $<10^{-10}$ & 0.71 & 0.58 \\
\bottomrule
\end{tabular}


\raggedright
\footnotesize
All KS tests are two-sided. Historical deterministic APS scores; a conservative split-conformal guarantee requires the held-out/exchangeability assumptions absent from the supplied protocol. For catastrophic tasks, test scores reportedly shift rightward; two-sided KS and mean shifts alone do not establish stochastic dominance or verify Theorem~\ref{thm:score_inflation}. Two robust tasks (i-incoterms, i-shippoint) show test scores shifted \textit{leftward}---the model becomes more confident during COVID, possibly due to reduced product diversity.
\end{table}

\section{RAPS Comparison}\label{app:raps}

Table~\ref{tab:raps} compares APS and RAPS coverage drops across all 8 tasks with 10-seed means ($\lambda=0.01$, $k_{\text{reg}}=5$).

\begin{table}[htbp]
\centering
\caption{APS vs.\ RAPS Coverage Drop (10-Seed Means $\pm$ Std). $\lambda=0.01$, $k_{\text{reg}}=5$.}\label{tab:raps}
\small
\begin{tabular}{@{}lccccc@{}}
\toprule
Task & Cl & APS Drop & RAPS Drop & $\Delta$ & Note \\
\midrule
\multicolumn{6}{l}{\textit{Severe tasks}} \\
s-shipcond & 45 & $60.4\pm31.8$ & $67.8\pm25.4$ & +7.4 & No improvement \\
s-group & 459 & $73.5\pm33.9$ & $10.4\pm1.3$ & $-$63.1 & RAPS protects \\
s-payterms & 137 & $79.9\pm29.4$ & $35.1\pm28.6$ & $-$44.8 & Partial recovery \\
\midrule
\multicolumn{6}{l}{\textit{Non-severe tasks}} \\
i-plant & 35 & $11.8\pm9.0$ & $13.1\pm5.8$ & +1.3 & Near-identical \\
i-shippoint & 69 & $9.3\pm27.6$ & $20.5\pm21.6$ & +11.2 & RAPS worse \\
i-incoterms & 13 & $10.8\pm9.1$ & $10.8\pm9.1$ & 0.0 & Identical \\
s-incoterms & 13 & $6.7\pm6.8$ & $6.7\pm6.8$ & 0.0 & Identical \\
s-office & 25 & $0.04\pm0.03$ & $0.04\pm0.03$ & 0.0 & Identical \\
\bottomrule
\end{tabular}


\raggedright
\footnotesize
All values are 10-seed means $\pm$ std. APS drops differ from 50-seed Table~\ref{tab:main_results} due to smaller seed range (10 vs.\ 50) and high variance. The key comparison is APS vs.\ RAPS within the same seed set. Note: RAPS experiments used a random 50/50 calibration split rather than the deterministic first-half/second-half split used in the main APS experiments (Appendix~\ref{app:reproducibility}); this protocol difference does not affect the within-experiment APS vs.\ RAPS comparison but means RAPS coverage values are not directly comparable to Table~\ref{tab:main_results}. RAPS eliminates the catastrophic drop for sales-group by truncating the class accumulation that amplifies score inflation in many-class settings (mean set size: 232/459 classes). Sales-shipcond is unhelped because single-feature dependence places dominant mass on the wrong class \textit{before} the $k_{\text{reg}}=5$ penalty activates---the exchangeability violation occurs at the top of the APS ordering, not through class accumulation. For low-class tasks ($\leq$35 classes), RAPS and APS produce identical or near-identical drops, confirming RAPS regularization has negligible effect when prediction sets are already small. Item-shippoint ($K=69$) is an exception: RAPS worsens coverage by 11.2~pp (paired Wilcoxon $p=0.064$, 10 seeds), reflecting that the regularization penalty can disrupt calibration for intermediate-cardinality tasks with high seed variance.
\end{table}

\section{Model Class Sensitivity}\label{app:rf}

To test whether SHAP concentration is diagnostic across model classes, we replicate the analysis with \texttt{RandomForestClassifier} (500 trees), \texttt{XGBClassifier}, and \texttt{CatBoostClassifier} (ordered boosting), all with seed=42 and the same data pipeline (Table~\ref{tab:rf}).

\begin{table}[htbp]
\centering
\caption{SHAP Concentration and Coverage Drop Across Model Classes.}\label{tab:rf}
\small
\begin{tabular}{@{}lcccccccc@{}}
\toprule
Task & LGB & RF & XGB & CB & LGB & RF & XGB & CB \\
 & \multicolumn{4}{c}{Concentration (\%)} & \multicolumn{4}{c}{Coverage Drop (\%)} \\
\midrule
s-shipcond & 50.7 & 57.4 & 44.6 & 53.7 & 71.6 & 15.8 & 47.2 & 1.1 \\
s-group & 47.3 & 22.5 & 33.3 & 31.8 & 71.2 & 50.0 & 56.8 & 14.5 \\
s-payterms & 54.2 & 27.9 & 30.6 & 25.7 & 77.1 & 13.8 & 20.4 & $-$5.3 \\
i-shippoint & 48.8 & 52.1 & 41.0 & 43.1 & 18.5 & 22.7 & 80.4 & 61.8 \\
i-plant & 23.9 & 26.8 & 26.0 & 29.0 & 10.6 & 15.3 & 9.7 & 8.5 \\
s-incoterms & 23.7 & 26.8 & 21.6 & 28.2 & 8.5 & 2.5 & 3.7 & $-$1.5 \\
i-incoterms & 28.9 & 17.5 & 18.4 & 21.7 & 11.3 & 2.9 & 1.1 & $-$0.4 \\
s-office & 42.6 & 24.5 & 36.8 & 33.5 & 0.1 & 0.0 & 0.0 & $-$0.0 \\
\bottomrule
\end{tabular}
\end{table}

\textbf{Results.} The diagnostic's predictive power varies by model class (Table~\ref{tab:model_sensitivity}). All three gradient-boosted implementations show positive correlation: LightGBM ($\rho = 0.833$, $p = 0.010$), CatBoost ($\rho = 0.667$, $p = 0.071$), and XGBoost ($\rho = 0.548$, $p = 0.16$). RandomForest ($\rho = 0.30$, $p = 0.47$) does not replicate.

\textbf{Mixed-effects analysis (historical, unverified).} The archived summary fits task random intercepts to 24 task/model combinations from eight tasks and reports $\beta_1=1.64$, Wald 95\% CI $[0.70,2.58]$, $p=0.0006$. With eight clusters, the adequacy of the asymptotic interval/test is uncertain. A previously suggested Kenward--Roger $p$ range was not computed and is withdrawn. The confidence interval and $p$-value share assumptions; one cannot claim the interval excludes zero under an uncomputed alternative correction. The original fitting code, model diagnostics and task/model data are missing. Random intercepts alone also do not establish independence of tasks sharing the benchmark.

\begin{table}[htbp]
\centering
\caption{Concentration--Drop Relationship by Model Class.}\label{tab:model_sensitivity}
\small
\begin{tabular}{@{}lccc@{}}
\toprule
Model & Method & Conc.--Drop $\rho$ & $p$ \\
\midrule
LightGBM & Boosting (GOSS) & \textbf{0.833} & \textbf{0.010} \\
CatBoost & Boosting (ordered) & 0.667 & 0.071 \\
XGBoost & Boosting (histogram) & 0.548 & 0.16 \\
RandomForest & Bagging & 0.30 & 0.47 \\
\midrule
\multicolumn{2}{@{}l}{Mixed-effects, 3 boosting ($n=24$)} & $\beta_1=1.64$ & \textbf{$<$0.001} \\
\multicolumn{2}{@{}l}{Mixed-effects, all 4 ($n=32$)} & $\beta_1=1.06$ & \textbf{0.010} \\
\bottomrule
\end{tabular}
\end{table}

\textbf{Interpretation.} The historical model/task summaries show heterogeneous outcomes and attribution profiles. They do not establish why RF or MLP correlations differ, nor that differences between correlations are statistically significant. Architectural explanations involving bagging smoothness or boosting reinforcement are hypotheses, requiring controlled comparisons and common attribution definitions.

The diagnostic is \textit{model-specific}: it measures how a particular model's learned dependence structure creates vulnerability. Practitioners should compute concentration for their \textit{deployed} model class rather than transferring results across model families.

\textbf{Neural network extension (historical, unverified).} The reported two-layer MLP uses permutation importance rather than TreeSHAP, with 10 seeds for sales tasks and 3 for item tasks. Its concentration ratio is therefore not the same attribution estimand as the tree-model SHAP ratio. The reported $\rho=0.43$, $p=0.29$ cannot isolate architecture from attribution-method differences. High-drop low-concentration examples motivate a possible multi-feature mechanism but do not prove it; original code/results are required to assess the claim.

\subsection{MLP vs LightGBM Task-Level Comparison}\label{app:mlp_lgb_comparison}

Table~\ref{tab:mlp_lgb_comparison} provides a structured comparison of SHAP concentration and coverage drop between MLP and LightGBM across all 8 SALT tasks.

\begin{table}[htbp]
\centering
\caption{Task-Level Comparison: MLP vs LightGBM. MLP\_C = MLP SHAP concentration; MLP\_Drop = MLP coverage drop (pp); LGB\_C = LightGBM SHAP concentration; LGB\_Drop = LightGBM coverage drop (pp); Mode = failure mode classification.}\label{tab:mlp_lgb_comparison}
\small
\begin{tabular}{@{}lccccc@{}}
\toprule
Task & MLP\_C & MLP\_Drop & LGB\_C & LGB\_Drop & Mode \\
\midrule
s-group & 27.5\% & 78.4pp & 47.3\% & 71.2pp & Global \\
s-payterms & 28.0\% & 60.6pp & 54.2\% & 77.1pp & Global \\
i-shippoint & 77.2\% & 82.3pp & 48.8\% & 18.5pp & Conc \\
s-incoterms & 30.3\% & 42.3pp & 23.7\% & 8.5pp & Mixed \\
s-shipcond & 53.1\% & 67.1pp & 50.7\% & 71.6pp & Conc \\
i-incoterms & 31.0\% & 24.1pp & 28.9\% & 11.3pp & Mixed \\
i-plant & 28.2\% & 26.7pp & 23.9\% & 10.6pp & Mixed \\
s-office & 0.0\% & 0.1pp & 42.6\% & 0.1pp & N/A \\
\bottomrule
\end{tabular}


\raggedright
\footnotesize
Mode: Conc = concentrated-dependence failure (C $>$ 40\% predicts drop); Global = global-sensitivity failure (C $\leq$ 40\% yet catastrophic drop); Mixed = intermediate regime; N/A = no meaningful drop in either model. Critical observation: s-group and s-payterms show MLP C $\approx$ 28\% yet 60--78pp drops, demonstrating the global-sensitivity failure mode where SHAP concentration does not capture the MLP's vulnerability mechanism.
\end{table}

\subsection{Failure Mode Definitions and Quantitative Criteria}\label{app:failure_modes}

We formalize two distinct failure modes that determine when SHAP concentration is diagnostic.

\smallskip
\noindent\fbox{\parbox{0.97\columnwidth}{%
\textbf{Definition 1: Concentrated-Dependence Failure}\\[2pt]
\textit{Conditions}: SHAP concentration $C > 40\%$; dominant feature undergoes distribution shift.\\
\textit{Diagnostic status}: \textbf{Applies}---SHAP concentration correctly predicts vulnerability.\\
\textit{Characteristic}: Model depends on a single feature whose OOD shift causes coverage collapse.
}}

\smallskip
\noindent\fbox{\parbox{0.97\columnwidth}{%
\textbf{Definition 2: Global-Sensitivity Failure}\\[2pt]
\textit{Conditions}: SHAP concentration $C \leq 40\%$; coverage degrades via coordinated multi-feature sensitivity.\\
\textit{Diagnostic status}: \textbf{Does NOT apply}---SHAP concentration fails to predict vulnerability.\\
\textit{Characteristic}: Model distributes importance but fails through global sensitivity to shift; requires alternative diagnostics (Table~\ref{tab:alternative_diagnostics}).
}}
\smallskip

\begin{table}[htbp]
\centering
\caption{Quantitative Criteria for Failure Mode Classification.}\label{tab:failure_mode_criteria}
\small
\begin{tabular}{@{}lcc@{}}
\toprule
Criterion & Concentrated & Global \\
\midrule
Top-1 SHAP & $C > 40\%$ & $C \leq 40\%$ \\
Gini coefficient & $> 0.6$ & $\leq 0.6$ \\
Top-3 cumulative & $> 70\%$ & $\leq 70\%$ \\
\bottomrule
\end{tabular}


\raggedright
\footnotesize
Concentrated failure mode: SHAP concentration diagnostic applies. Global failure mode: alternative diagnostics required (see Table~\ref{tab:alternative_diagnostics}).
\end{table}

\begin{table}[htbp]
\centering
\caption{Decision Rule: When Does the SHAP Concentration Diagnostic Apply?}\label{tab:decision_rule}
\small
\begin{tabular}{@{}llll@{}}
\toprule
Model Class & Condition & Prediction & Recommendation \\
\midrule
\multirow{2}{*}{\shortstack[l]{GBM/XGBoost/\\LightGBM/CatBoost}} & $C > 40\%$ & \textbf{VULNERABLE} & Mitigate or monitor \\
 & $C \leq 40\%$ & Lower risk & Standard deployment \\
\midrule
\multirow{2}{*}{MLP/Neural Net} & $C > 40\%$ & VULNERABLE (likely) & Mitigate or monitor \\
 & $C \leq 40\%$ & \textit{Inconclusive} & Use gradient-based methods \\
\bottomrule
\end{tabular}


\raggedright
\footnotesize
The diagnostic is validated for gradient-boosted classifiers ($\rho=0.833$, $p=0.010$ for LightGBM). For MLPs, low concentration does not ensure robustness due to global-sensitivity failures ($\rho=0.43$, $p=0.29$).
\end{table}

\subsection{Alternative Diagnostics for MLPs When $C \leq 40\%$}\label{app:alternative_diagnostics}

When SHAP concentration indicates low risk ($C \leq 40\%$) but the model is an MLP or other neural network, the following alternative diagnostics may capture global-sensitivity failures (Table~\ref{tab:alternative_diagnostics}).

\begin{table}[htbp]
\centering
\caption{Alternative Diagnostics for Global-Sensitivity Failures in Neural Networks.}\label{tab:alternative_diagnostics}
\small
\begin{tabular}{@{}lp{8cm}@{}}
\toprule
Diagnostic & Description \\
\midrule
Gradient sensitivity & $S_j = \mathbb{E}[\|\partial f / \partial x_j\|_2]$: average gradient magnitude per feature \\
Feature ablation & $\Delta_j = \text{Cov}(D) - \text{Cov}(D|x_j=\mu_j)$: coverage change when ablating feature $j$ \\
Lipschitz bound & $\text{Cov}_{\text{test}} \geq \text{Cov}_{\text{val}} - L \cdot W_2(D_{\text{val}}, D_{\text{test}})$: coverage lower bound via Lipschitz constant $L$ and Wasserstein-2 distance \\
\bottomrule
\end{tabular}


\raggedright
\footnotesize
These diagnostics target global-sensitivity failures where coordinated multi-feature shifts drive coverage degradation. They are proposed as complements to SHAP concentration for neural network deployments; empirical validation is left to future work.
\end{table}

\section{Hyperparameter Sensitivity Analysis}\label{app:hp_sensitivity}

To assess whether the SHAP concentration--coverage drop correlation is robust to hyperparameter choices, we evaluate four LightGBM configurations on the 8 SALT multiclass tasks (Table~\ref{tab:hp_sensitivity}).

\begin{table}[htbp]
\centering
\caption{Hyperparameter Sensitivity Analysis (SALT, $n=8$ tasks). All configurations use boosting=gbdt, feature\_fraction=0.8, bagging\_fraction=0.8, bagging\_freq=5.}\label{tab:hp_sensitivity}
\small
\begin{tabular}{@{}lccccccc@{}}
\toprule
Config & leaves & lr & n\_est & $\rho$ & $p$ & Thresh & F1 \\
\midrule
Default & 31 & 0.05 & 100 & 0.833 & 0.010 & 45.0\% & 1.0 \\
Deeper & 63 & 0.05 & 100 & 0.810 & 0.015 & 47.5\% & 1.0 \\
Faster & 31 & 0.10 & 100 & 0.810 & 0.015 & 45.0\% & 1.0 \\
More & 31 & 0.05 & 200 & 0.833 & 0.010 & 45.0\% & 1.0 \\
\bottomrule
\end{tabular}


\raggedright
\footnotesize
leaves = num\_leaves; lr = learning\_rate; n\_est = num\_boost\_round; Thresh = optimal threshold for at-risk classification (drop $>$15~pp); F1 = F1 score at that threshold.
\end{table}

\textbf{Results.} Across all four hyperparameter configurations, the concentration--drop correlation remains strong and significant: $\rho = 0.821 \pm 0.012$ (mean $\pm$ std), with $p \leq 0.015$ for all configurations. The optimal threshold varies minimally ($45.6\% \pm 1.1\%$), and all configurations achieve perfect F1 = 1.0 for at-risk classification on SALT tasks.

\textbf{Interpretation.} The four reported hyperparameter settings provide a small retrospective sensitivity summary. They do not establish a structural property of gradient boosting or robustness to unrestricted tuning. The generating artifacts are missing, and the ablation default uses 100 boosting rounds versus the main experiment's stated maximum of 500; those protocols require reconciliation.

\subsection{Per-Task Concentration Stability Across HP Configurations}\label{app:pertask_hp_stability}

Table~\ref{tab:pertask_hp_concentration} shows per-task SHAP concentration values across all four hyperparameter configurations.

\begin{table}[htbp]
\centering
\caption{Per-Task SHAP Concentration Stability Across HP Configurations. Std = standard deviation across four configurations.}\label{tab:pertask_hp_concentration}
\small
\begin{tabular}{@{}lccccc@{}}
\toprule
Task & Default & Deeper & Faster & More & Std \\
\midrule
sales-payterms & 54.2\% & 57.3\% & 55.9\% & 54.3\% & 1.3\% \\
sales-shipcond & 50.7\% & 53.7\% & 52.3\% & 48.5\% & 1.9\% \\
item-shippoint & 48.8\% & 50.0\% & 50.7\% & 46.0\% & 1.8\% \\
sales-group & 47.3\% & 49.4\% & 48.5\% & 45.2\% & 1.6\% \\
sales-office & 42.6\% & 45.3\% & 43.0\% & 42.6\% & 1.1\% \\
item-incoterms & 28.9\% & 30.1\% & 28.9\% & 29.3\% & 0.5\% \\
item-plant & 23.9\% & 25.2\% & 24.4\% & 23.9\% & 0.5\% \\
sales-incoterms & 23.7\% & 25.4\% & 25.2\% & 22.5\% & 1.2\% \\
\bottomrule
\end{tabular}


\raggedright
\footnotesize
Tasks sorted by Default concentration. Default = num\_leaves=31, lr=0.05, n\_est=100; Deeper = num\_leaves=63; Faster = lr=0.10; More = n\_est=200. Mean within-task std = 1.2\%, demonstrating high stability of concentration values across HP configurations. The rank order of tasks by concentration is preserved across all four configurations.
\end{table}

\section{Shift Characterization}\label{app:shift_characterization}

To characterize the nature of distribution shift in SALT tasks, we decompose shift into covariate and concept components using standard metrics (Table~\ref{tab:shift_characterization}).

\begin{table}[htbp]
\centering
\caption{Shift Characterization for SALT Tasks. Cov KS = mean Kolmogorov--Smirnov statistic across features (train vs.\ test); \%Sig = fraction of features with KS $p < 0.05$; Label JS = Jensen--Shannon divergence of label distributions; Acc Drop = test accuracy minus validation accuracy.}\label{tab:shift_characterization}
\small
\begin{tabular}{@{}lccccc@{}}
\toprule
Task & Cov KS & \%Sig & Label JS & Acc Drop & Classification \\
\midrule
s-shipcond & 0.50 & 50\% & 0.083 & +11.9\% & Cov + Concept \\
s-group & 0.50 & 50\% & 0.330 & +1.0\% & Cov (dominant) \\
s-payterms & 0.50 & 50\% & 0.078 & +4.0\% & Cov (dominant) \\
i-plant & 0.86 & 86\% & 0.144 & $-$8.7\% & Cov (dominant) \\
i-shippoint & 0.86 & 86\% & 0.152 & +13.1\% & Cov + Concept \\
s-incoterms & 0.50 & 50\% & 0.085 & +7.6\% & Cov + Concept \\
i-incoterms & 0.86 & 86\% & 0.140 & +63.4\% & Cov + Concept \\
s-office & 0.50 & 50\% & 0.026 & $-$0.2\% & Cov (dominant) \\
\bottomrule
\end{tabular}


\raggedright
\footnotesize
Cov = covariate shift; Concept = concept shift (change in $P(Y|X)$). Accuracy changes can be positive under shift when the test distribution favors majority classes.
\end{table}

\textbf{Key findings.}
\begin{itemize}[nosep,leftmargin=*]
    \item \textbf{Covariate shift is universal}: All SALT tasks exhibit substantial covariate shift (KS $\geq 0.50$, $\geq$50\% of features significantly different). This is expected given the COVID-19 temporal split.
    \item \textbf{Concept shift varies}: Label distribution shift (JS divergence) ranges from 0.026 (s-office) to 0.330 (s-group). Tasks with both high covariate shift and substantial label shift (s-shipcond, i-incoterms) tend toward catastrophic failure.
    \item \textbf{Accuracy is not diagnostic}: Accuracy changes are uninformative for coverage prediction---s-shipcond shows +11.9\% accuracy improvement yet 71.6~pp coverage drop, while s-office shows $-$0.2\% accuracy change and 0.1~pp coverage drop. This occurs because conformal prediction set size is sensitive to calibration--test score distribution mismatch, which accuracy does not capture.
\end{itemize}

\textbf{Interpretation.} Standard shift detection identifies covariate shift uniformly across all tasks (consistent with Section~\ref{sec:shift_detection}), while concept shift magnitude partially distinguishes catastrophic from robust outcomes. However, neither metric alone predicts failure severity as effectively as SHAP concentration ($\rho = 0.833$), which captures how the model's learned dependence structure interacts with the specific features that shift. The worst outcomes occur when concentrated dependence targets a feature experiencing both covariate and concept shift.

\subsection{Per-Feature Covariate Shift Details}\label{app:perfeature_shift}

Tables~\ref{tab:sales_feature_shift} and~\ref{tab:item_feature_shift} show per-feature KS statistics for sales-level and item-level tasks respectively.

\begin{table}[htbp]
\centering
\caption{Per-Feature Covariate Shift for Sales-Level Tasks (6 features). KS = Kolmogorov--Smirnov statistic (train vs.\ test).}\label{tab:sales_feature_shift}
\small
\begin{tabular}{@{}lcc@{}}
\toprule
Feature & Shifted & KS \\
\midrule
SALESDOCUMENTTYPE & No & 0.0 \\
SALESORGANIZATION & \textbf{Yes} & 1.0 \\
DISTRIBUTIONCHANNEL & \textbf{Yes} & 1.0 \\
ORGANIZATIONDIVISION & \textbf{Yes} & 1.0 \\
BILLINGCOMPANYCODE & No & 0.0 \\
TRANSACTIONCURRENCY & No & 0.0 \\
\bottomrule
\end{tabular}


\raggedright
\footnotesize
KS = 1.0 indicates complete distribution non-overlap between train and test. 3/6 features (50\%) show significant covariate shift, consistent with the \%Sig values in Table~\ref{tab:shift_characterization}.
\end{table}

\begin{table}[htbp]
\centering
\caption{Per-Feature Covariate Shift for Item-Level Tasks (7 features). KS = Kolmogorov--Smirnov statistic (train vs.\ test).}\label{tab:item_feature_shift}
\small
\begin{tabular}{@{}lcc@{}}
\toprule
Feature & Shifted & KS \\
\midrule
SALESDOCUMENTITEM & \textbf{Yes} & 1.0 \\
SALESDOCUMENTITEMCATEGORY & No & 0.0 \\
PRODUCT & \textbf{Yes} & 1.0 \\
SOLDTOPARTY & \textbf{Yes} & 1.0 \\
SHIPTOPARTY & \textbf{Yes} & 1.0 \\
BILLTOPARTY & \textbf{Yes} & 1.0 \\
PAYERPARTY & \textbf{Yes} & 1.0 \\
\bottomrule
\end{tabular}


\raggedright
\footnotesize
6/7 features (86\%) show significant covariate shift, consistent with the \%Sig values in Table~\ref{tab:shift_characterization}. The high proportion of shifted features at item-level reflects COVID-19's impact on supply chain transactions: new products, parties, and item identifiers entered the system during the pandemic period.
\end{table}

\section{WILDS Benchmark Validation}\label{app:wilds}

To further validate the specificity of the SHAP concentration diagnostic, we evaluate on four additional benchmark datasets from the WILDS benchmark~\citep{koh2021wilds} and standard repositories. The historical summaries report low concentration and robust coverage in four evaluations. Their generating code and outputs are absent; Covertype is reused under another split, and CivilComments/Adult are binary tasks outside the selected primary subset.

\begin{table}[htbp]
\centering
\caption{WILDS and Standard Benchmark Validation (Specificity Test). Two evaluations use WILDS; two use standard datasets. Covertype also appears in the primary analysis under a different split. All datasets exhibit low SHAP concentration (\textbf{all $<$40\%}; mean = 17.0\%) and robust coverage, reporting four true negatives among this small selected sample; independent specificity is not established. Shift types: Demographic = identity-based subgroup shift; User/Time = user and temporal domain shift; Geographic = regional covariates; Age = demographic age-based shift. Note: This table tests specificity (low-$C$ $\rightarrow$ robust); for the catastrophic Covertype case (high-$C$ $\rightarrow$ failure), see Table~\ref{tab:framework_validation}.}\label{tab:wilds_validation}
\small
\begin{tabular}{@{}lcccccc@{}}
\toprule
Dataset & Source & Shift Type & Conc.\ (\%) & Drop (pp) & Predicted & Actual \\
\midrule
CivilComments & WILDS & Demographic & 26.2 & $-$0.7 & Robust & Robust \\
Amazon & WILDS & User/Time & 1.9 & +0.4 & Robust & Robust \\
Covertype$^\dagger$ & sklearn & Geographic & 12.5 & +8.0 & Robust & Robust \\
Adult & OpenML & Age-based & 27.5 & +1.4 & Robust & Robust \\
\bottomrule
\end{tabular}

\raggedright
\footnotesize
$^\dagger$Covertype here uses a different split protocol from the main external validation (Table~\ref{tab:external_datasets}), which explains the different concentration and drop values. The main-text Covertype uses wilderness areas 1--2/3/4 geographic split (concentration 49.8\%, drop 81.8~pp); this validation uses a wilderness area 1 vs.\ others geographic split (concentration 12.5\%, drop 8.0~pp). Both results are consistent with the diagnostic: high concentration predicts failure, low concentration predicts robustness.
\end{table}

\textbf{Dataset details.}
\begin{itemize}[nosep,leftmargin=*]
    \item \textbf{CivilComments} (WILDS): Toxicity classification with demographic identity-based shift. TF-IDF features (1,000 dimensions). Low concentration (26.2\%) correctly predicts maintained coverage.
    \item \textbf{Amazon Reviews} (WILDS): Product rating prediction with user/time domain shift. TF-IDF features (1,000 dimensions). Very low concentration (1.9\%) correctly predicts robustness.
    \item \textbf{Covertype} (sklearn): Forest cover classification with geographic shift (wilderness area 1 vs.\ others). 54 cartographic features. Low concentration (12.5\%) correctly predicts mild degradation.
    \item \textbf{Adult Income} (OpenML): Income prediction with age-based demographic shift (train: age$<$50, test: age$\geq$50). 14 demographic/employment features. Low concentration (27.5\%) correctly predicts maintained coverage.
\end{itemize}

\textbf{Specificity validation.} The four reported benchmark evaluations have low SHAP concentration---\textbf{all below the 40\% threshold} (individual values: 26.2\%, 1.9\%, 12.5\%, 27.5\%; \textbf{mean concentration = 17.0\%}). The diagnostic correctly predicts robustness for all four, giving an observed 4/4 true-negative rate, with substantial sampling and selection uncertainty and without raw-run verification. This complements the sensitivity evidence from SALT, where high-concentration tasks correctly predict catastrophic failure.

\textbf{Combined validation accuracy (SALT + WILDS).} Combining the 8 SALT tasks with these 4 WILDS benchmarks yields $n=12$ total tasks (distinct from the $n=16$ primary endpoint which uses 8 different external UCI datasets). The threshold accuracy is \textbf{91.7\% (11/12)}:
\begin{itemize}[nosep,leftmargin=*]
    \item \textbf{Sensitivity} (SALT): 4/4 at-risk SALT tasks flagged under the $>$15~pp rule (historical descriptive counts)
    \item \textbf{Specificity} (WILDS + robust SALT): 7/8 robust tasks correctly unflagged (87.5\%)
    \item \textbf{Single misclassification}: sales-office (concentration 42.6\%, protective secondary feature)
\end{itemize}

These small, selected historical counts do not validate an operational boundary. Missing runs, binary tasks outside the primary scope and reuse of Covertype require explicit treatment in a prospective evaluation.

\end{document}
